\section{DQN 用慢目标与回放缓和训练耦合}
\label{sec:16-Control-and-Reinforcement-Learning/02-Reinforcement-Learning/07}

换成在线的场景：一个像素输入的游戏环境，你接了个卷积网络输出所有动作的 \(Q\) 值，按环境时间顺序每走一步做一次梯度更新，\(\varepsilon\)-贪心探索。跑了两小时，分数先升后崩；Q 值曲线一节一节往上抬，网络在几十步之内反复被同一小段走廊的画面刷屏。你调小学习率，崩溃推迟；你换个种子，崩溃换个时刻。上一节已经说过，三角齐全时期望更新本身就可能不收敛，所以这里没什么好意外的。

意外的是，加两个零件之后同一份代码能跑通。DQN 的两个零件不动三角的任何一角，值得弄清楚它们各自动了什么。

\subsection{先把训练里的耦合数清楚}

对一条经验 \(e_t=(S_t,A_t,R_{t+1},S_{t+1},D_{t+1})\)，其中 \(D_{t+1}\) 标记是否终止，回归目标是

\[
Y_t=R_{t+1}+\gamma(1-D_{t+1})\max_{a'}Q(S_{t+1},a';\theta^-),
\]

在线参数拟合当前动作那一路输出：

\[
L(\theta)=\E_{e\sim\mathcal B}\big[\ell\big(Y(e;\theta^-)-Q(S,A;\theta)\big)\big],
\]

其中 \(\mathcal B\) 是回放缓冲的采样分布，\(\theta^-\) 是目标网络参数。

这个式子里藏着三条耦合。标签 \(Y\) 由参数算出，参数一动标签就动，这是自举那一角。采样分布 \(\mathcal B\) 由行为策略决定，而 \(\max_{a'}\) 评估的是贪心策略，这是离策略那一角。同一批参数还决定下一步采集什么数据，数据又反过来决定参数怎么动，这条闭环是强化学习独有的、监督学习里没有的。三条耦合的时间尺度原本都等于一次梯度更新，DQN 做的事情是把其中两条拉长。

\subsection{目标网络把标签在一段时间内钉住}

令 \(\theta^-=\theta\) 时，一次梯度步同时移动预测和标签：你朝目标走一步，目标也跟着挪一点。硬同步每 \(C\) 步做一次 \(\theta^-\leftarrow\theta\)，软同步用

\[
\theta^-\leftarrow(1-\tau)\theta^-+\tau\theta,\qquad 0<\tau\ll1,
\]

两种写法给出同一件东西：在若干次更新之内，标签近似固定，优化器面对的是一个普通的回归问题，梯度下降的常规直觉重新可用（\hyperref[sec:09-Convex-Optimization/06-Optimization-Algorithms/01]{``梯度下降真正难在步长''}）。同步周期 \(C\) 或 \(\tau\) 是一个时间尺度旋钮：太快，标签重新变成移动靶；太慢，网络在长时间里拟合一个陈旧的目标，奖励信息传播得很吃力。

实现上必须切断 \(Y_t\) 对 \(\theta^-\) 的梯度。忘了 \texttt{detach}，优化器会发现同时压低标签和抬高预测更省力，跑的算法就不再是上面写的那个了。

\begin{center}
\begin{tikzpicture}[>=stealth,scale=0.85]
% 左：目标冻住 vs 目标跟着跑
\draw[->] (0,0) -- (5.4,0) node[below,font=\small] {更新步};
\draw[->] (0,0) -- (0,3.4) node[left,font=\small] {\(Q\)};
\draw[thick] (0,1.0) -- (1.25,1.0) (1.25,1.5) -- (2.5,1.5)
  (2.5,1.8) -- (3.75,1.8) (3.75,1.95) -- (5.0,1.95);
\draw[dashed] (1.25,1.0) -- (1.25,1.5) (2.5,1.5) -- (2.5,1.8)
  (3.75,1.8) -- (3.75,1.95);
\draw[thick] plot[smooth] coordinates
  {(0,0.2) (0.4,0.66) (0.8,0.9) (1.25,0.98) (1.7,1.28) (2.1,1.44)
   (2.5,1.49) (2.9,1.69) (3.4,1.78) (3.75,1.79) (4.2,1.9) (4.7,1.94) (5.0,1.95)};
\draw[dotted,thick] plot[domain=0:4.6,samples=40] (\x,{0.22*exp(0.58*\x)});
\node[font=\small,align=left] at (2.4,3.05) {\(\theta^-=\theta\)：标签一起跑};
\draw[->] (3.6,3.0) -- (4.2,2.65);
\node[font=\small] at (4.2,1.3) {冻住的标签};
% 右：回放打散相关
\begin{scope}[shift={(7.0,0.2)}]
\draw (0,2.1) rectangle (4.8,2.7);
\foreach \x in {0.2,0.6,1.0,1.4,1.8,2.2,2.6,3.0,3.4,3.8,4.2,4.6}
  \draw (\x,2.1) -- (\x,2.7);
\fill[gray!45] (0.6,2.1) rectangle (1.4,2.7);
\node[below,font=\small,align=center] at (2.4,2.05) {按时间顺序取 batch：一小段轨迹主导梯度};
\draw (0,0.7) rectangle (4.8,1.3);
\foreach \x in {0.2,0.6,1.0,1.4,1.8,2.2,2.6,3.0,3.4,3.8,4.2,4.6}
  \draw (\x,0.7) -- (\x,1.3);
\foreach \x in {0.2,1.4,2.2,3.4,4.2}
  \fill[gray!45] (\x,0.7) rectangle (\x+0.4,1.3);
\node[below,font=\small,align=center] at (2.4,0.65) {均匀回放取 batch：横跨许多旧策略};
\end{scope}
\end{tikzpicture}
\end{center}

\emph{左边把标签的移动放慢，右边把样本的相关性打散；两处改的都是耦合的强度与节奏。}

\subsection{回放把一段轨迹换成一个混合分布}

按时间顺序更新时，一个 batch 里的画面几乎是同一秒的同一个走廊，梯度反复指向同一个方向。缓冲区保存最近若干万条转移再随机抽样，图右上那条带子于是变成图右下那条：一个 batch 的样本来自缓冲区里的许多时刻，也来自许多个已经作废的旧策略。相关性下降，一次昂贵交互还能被用上很多次。

麻烦同时产生。缓冲区里混着的旧策略让算法天然是离策略的，容量越大，采样分布离当前策略越远；容量太小又会很快忘掉稀有经验。均匀回放没有把数据变回独立同分布，也不保证覆盖贪心策略需要访问的状态动作对，它只是削弱了局部时间相关。优先回放按 \(\abs{\delta}\) 的大小抽样，把算力集中在当前误差大的样本上，采样分布则被进一步扭曲，需要用基于采样概率的重要性权重把偏差校回来——这套权重和下一节要讲的是同一件工具。大 TD 误差也未必意味着更值得学，噪声、异常值和标错的奖励同样会产生大误差。

\subsection{两个零件都没有拆掉三角}

回到上一节那张三角图：目标网络仍然自举，回放仍然离策略，卷积网络仍然是函数逼近。三个角一个没少。目标网络做的是把自举那一角的更新频率降下来，让标签的漂移慢于网络的拟合；回放做的是把采样那一角摊平，让梯度不再被一小段轨迹绑架。耦合被削弱，反馈回路的增益被压低，误差绕圈放大的速度可能慢到崩溃发生在训练结束之后。

这是工程缓解，没有收敛定理跟在后面。它在什么条件下够用，其实可以对账：在线交互不断把新策略的数据刷进缓冲区，覆盖随策略改进而更新，贪心动作大体落在数据支持内，这时慢下来的耦合确实能撑住。反过来，在一份固定的离线数据集上训练，条件全部失效——\(\max_{a'}\) 会挑中日志里从未出现的动作，网络在那里的外推值凭空虚高，虚高值再被自举一路往前传，而没有任何新样本能来纠正它。数据再多也补不出行为策略从未尝试过的反事实。离线场景需要别的办法：把新策略约束在数据支持附近，对分布外动作取保守估计，或者退回到行为克隆这样目标更明确的基线。

\subsection{Double DQN 拆开选择与评价}

标准目标里的 \(\max_{a'}Q(S',a';\theta^-)\) 有一个与三角无关的独立毛病。设各动作的估计含独立噪声，由 Jensen 不等式，

\[
\E\Big[\max_{a'}\widehat Q(S',a')\Big]\ge\max_{a'}\E\big[\widehat Q(S',a')\big],
\]

取最大值这个操作偏爱正误差大的动作，动作数越多、噪声越大，乐观偏差越大。Double DQN 用在线网络选动作、用目标网络评估它：

\[
a^*=\argmax_{a'}Q(S',a';\theta),
\qquad
Y^{\mathrm{Double}}=R+\gamma(1-D)\,Q(S',a^*;\theta^-).
\]

选择与评价不再共享同一份噪声，过估计通常明显减轻。两套网络仍然相关（一个是另一个的滑动平均），所以偏差不会归零，某些任务上还会翻成低估。把它记成针对最大化选择偏差的一个专门修正，别把它当成万能稳定器。

\subsection{损失函数与掩码管的是数值，不是方向}

平方损失下梯度幅度随 TD 误差线性增长，一个奖励尖峰或一次早期的错误标签就能让少数样本主导整批更新。Huber 损失

\[
\ell_\kappa(\delta)=
\begin{cases}
\tfrac12\delta^{2}, & \abs{\delta}\le\kappa,\\[2pt]
\kappa\big(\abs{\delta}-\tfrac12\kappa\big), & \abs{\delta}>\kappa
\end{cases}
\]

在零附近保持平方、在尾部转成线性，梯度范数被 \(\kappa\) 封顶。梯度裁剪、奖励缩放和合理初始化做的是同一类事：防止一次异常更新把参数打飞。它们不会把错的 Bellman 目标变对，而长期依赖裁剪还会掩盖价值尺度已经失控的事实。

终止掩码那个 \((1-D_{t+1})\) 也值得单独说一句。环境真终止时未来价值为零，掩码应当生效；训练器因为步数上限截断回合时，未来价值并不为零，掩码不该生效。把超时一律当作终止，会系统性地低估长视界任务的价值，而这种低估看上去和三角带来的问题完全不像。

\subsection{损失下降不等于策略变好}

回归损失是按缓冲区分布加权的。损失很小，可能只说明网络把缓冲区里常见的转移拟合好了，并不说明贪心动作在缓冲区没覆盖的状态上正确。反过来，策略变强时损失往往上升，因为新策略把智能体带进了还没拟合过的区域。

值得同时盯住的量有几个：

\begin{itemize}
\tightlist
\item 多个随机种子下回合回报的分布，而不只是最好那条曲线；
\item 网络输出的 \(Q\) 值与实际 Monte Carlo 回报是否同量级；
\item TD 误差的分位数，尾部比均值先出问题；
\item 在线网络与目标网络预测的差距，它是标签漂移速度的直接读数。
\end{itemize}

这些量能告诉你耦合是否被压住，却没有一个能告诉你行为策略与目标策略差了多远。DQN 从头到尾把这个差异留在原地，用``缓冲区里的数据大致够用''这一默认假设兜着。下一节把这个差异写成一个可以计算的概率比率，同时引入一种新的不稳定。
